\section{Sharp thresholds for all cuboids}\label{sec:cuboids}
For ordered side lengths $2\le a\le b\le c$, let $F(a,b,c)$ be the least
inventory size guaranteeing a frame, and $G(a,b,c)$ the corresponding
complete-box threshold. These extend the cube notation: $F(n,n,n)=F(n)$.
A frame has eight corners and the noncorner positions on its twelve edges,
so its cardinality is
\[
 L=4(a+b+c)-16.
\]
The target may depend on the inventory. Every target has exactly three
exterior colours, and internal faces remain unrestricted.

\begin{theorem}[Cuboid thresholds]\label{thm:cuboids}
For integers $2\le a\le b\le c$, put $K(c)=\lfloor(24c-1)/5\rfloor$. Then
\begin{align*}
 F(a,b,c)&=\max\{16a-3,\;4(a+b+c)-16,\;K(c)\},\\
 G(a,b,c)&=\max\{abc,\;16a-3,\;K(c)\}.
\end{align*}
\end{theorem}
The three frame bounds come from a colour bucket, the number of positions,
and a balanced checkerboard collection. The shortest and longest sides
control different obstructions. For the upper bound, a weighted graph
chooses directions for a $D$ target whenever such a corner solution exists.
Otherwise the same bucket constraints provide a $Q$ or $H$ target.
This proof uses the geometry of Section~\ref{sec:geometry}, including the bucket lemma,
and the edge facts in Section~\ref{sec:frames}.

\subsection{Three lower bounds}
Cardinality gives $F(a,b,c)\ge L$. Fix one ordinary orientation class and
take $4a-1$ copies of each of its four $T$ varieties. By
Lemma~\ref{lem:balance}, these $16a-4$ sources can only assemble a $Q$
target. Its eight corners and four monochromatic edges parallel to an
axis of length $d$ require $8+4(d-2)=4d\ge4a$ sources from one bucket.
That bucket has only $4a-1$ copies. Hence $F(a,b,c)\ge16a-3$.

For the third bound, distribute $N_0=K(c)-1$ sources as evenly as possible
among the six checkerboards. Each bichromatic type then has at most
\[
 N_0-\lfloor N_0/6\rfloor=\lceil5N_0/6\rceil\le4c-1
\]
suppliers, since $5N_0\le24c-6$. All source corners are trichromatic.
Thus every supported target must have equal-coloured opposite faces,
of type $D$: equal colours on adjacent faces would produce a bicoloured
corner. In every axis assignment, the eight corners and $4(c-2)$ positions
on the longest edges require one common bichromatic pair. Their demand
is $4c$, exceeding its supply. Therefore $F(a,b,c)\ge K(c)$.

\subsection{Extending an opposite-face corner solution}
Reduce the sources by corner containment. As in Section~\ref{sec:frames},
this preserves edge availability and lifts every assembly to distinct
original copies. Index the checkerboards by their missing colour pairs,
viewing their multiplicities $x_e$ as edge weights on $K_4$.
The incident weight $z_v=\sum_{e\ni v}x_e$ counts checkerboards supporting
both orientations of the triple omitting $v$.
If $A_0,A_1$ count the two ordinary classes, that triple has a $D$ corner
solution exactly when
\begin{equation}\label{eq:cuboid-good}
 A_0+z_v\ge4,\qquad A_1+z_v\ge4,\qquad A_0+A_1+z_v\ge8.
\end{equation}
These are Hall's three inequalities for the two corner types, each with
four occurrences. Call such a vertex $v$ \emph{good}; in particular,
$z_v\ge8$ makes $v$ good.

\begin{lemma}[Weighted triangle]\label{lem:weighted-triangle}
Let $N\ge\max\{29,L,K(c)\}$ and $s_d=N-4d$ for $d\in\{a,b,c\}$.
Give the edges of $K_4$ nonnegative integer weights with sum at most $N$.
Suppose a nonempty set of vertices is designated good and includes every
vertex of incident weight at least eight. Some good vertex has a
complementary triangle whose edge weights, in increasing order, are at
most $s_c,s_b,s_a$, respectively.
\end{lemma}
\begin{proof}
We first record the numerical bounds
\begin{equation}\label{eq:cuboid-slack}
\begin{gathered}
 s_c\ge5,\qquad N\le6s_c+5,\qquad
 s_b\ge\max\{7,2b-1\},\\
 s_a\ge4b-1,\qquad s_c\ge\lfloor(4b-1)/3\rfloor.
\end{gathered}
\end{equation}
For $c\le6$, the first follows from $N\ge29$; for $c\ge7$, use
$K(c)=4c+\lfloor(4c-1)/5\rfloor\ge4c+5$.
The inequality $5K(c)\ge24c-5$ gives the second.
For $b\le3$, $s_b\ge17$; for $b\ge4$,
$s_b\ge4(a+c)-16\ge4b-8\ge\max\{7,2b-1\}$.
If $c\ge4$, then $s_a\ge4(b+c)-16\ge4b$; otherwise
$s_a\ge17\ge4b-1$.
Finally, for $b\le3$ the last floor is at most three; for $b\ge4$,
$s_c\ge4(a+b)-16\ge4b-8\ge\lfloor(4b-1)/3\rfloor$.

Suppose first that every edge weight is at most $s_b$. If no good vertex
has a complementary edge of weight at most $s_c$, all three complementary
edges of any good vertex have weight at least $s_c+1$. Each other vertex
then has incident weight at least $2(s_c+1)\ge12$, so all four vertices
are good. Every edge is complementary to one of them, forcing total
weight at least $6s_c+6>N$. This contradiction gives the required light
edge; the other two already fit $s_b$ and $s_a$.

Otherwise an edge has weight at least $s_b+1\ge8$. Either endpoint $v$
is good. The other five edges have total weight at most $4b-1$.
Writing $y_1\le y_2\le y_3$ for the weights complementary to $v$, we have
\[
 y_1\le\lfloor(4b-1)/3\rfloor\le s_c,\qquad
 y_2\le2b-1\le s_b,\qquad y_3\le4b-1\le s_a.
\]
\end{proof}

\begin{corollary}\label{cor:cuboid-D}
If $N\ge\max\{29,L,K(c)\}$ reduced sources have any $D$ corner solution,
they have an $a\times b\times c$ frame.
\end{corollary}
\begin{proof}
Apply Lemma~\ref{lem:weighted-triangle} with goodness defined
by~\eqref{eq:cuboid-good}. Assign the complementary colour pairs to the
three edge directions according to the lemma. Every such assignment is
realised by a $D$ face colouring: permutations of the triangle's vertices
induce all permutations of its edges. This chooses a target; each source
is still placed by a proper rotation.

Reserve eight sources for its corners. Each possesses all three target
pairs, so the residual supply for the pair assigned to a direction of
length $d$ is at least
\[
 N-x_e-8\ge4(d-2).
\]
Any two distinct pair neighbourhoods together contain every residual
source, since no reduced source lacks two bichromatic types. The total
edge demand is $L-8\le N-8$. Thus Hall's theorem fills the remaining
positions. The chosen target and axes may differ from those of the
initial corner solution.
\end{proof}

\subsection{Reservations when no \texorpdfstring{$D$}{D} corner solution exists}
Take
\[
 N=\max\{16a-3,L,K(c)\},
\]
and suppose no $D$ corner solution exists. Let $A\ge B$ count the ordinary
classes and $C$ the checkerboards. Equation~\eqref{eq:cuboid-good} gives
$z_v\le7$, hence $C\le14$ and $A\ge8$. It then gives $B+z_v\le3$ for
every $v$. Consequently
\begin{equation}\label{eq:cuboid-minority}
 4B+2C\le12,\qquad M:=B+C\le6,\qquad A=N-M\ge N-6\ge23.
\end{equation}
Each individual checkerboard multiplicity is at most three.
After reserving $r$ sources, each bichromatic type has at least $N-r-3$
suppliers; any two distinct types together see all $N-r$ remaining
sources. Hall therefore completes any residual bichromatic requirements
whose individual demands are at most $N-r-3$ and whose total is at most
$N-r$.

Use distinct colours $u,v,w$ and put $m=a-2$, $t=b-2$, $q=c-2$.
Table~\ref{tab:cuboid-demands} gives two targets, with face tuples in the
usual order and axis lengths $(a,b,c)$.

\begin{table}[ht]
\centering
\small
\begin{tabular}{llll}
\toprule
Target and face tuple & Monochromatic demand & Bichromatic demands & $r$\\
\midrule
$Q:(v,w,u,u,u,u)$ & $uu:4m$ & $uv,uw:2t+2q$ & $4a$\\
$H:(v,w,u,u,u,v)$ & $uu:2m,\ vv:t$ &
$uv:2m+t+2q$ & $2a+b+2$\\
 & & $uw:t+2q,\ vw:t$ & \\
\bottomrule
\end{tabular}
\caption{Noncorner cuboid demands. The reservation $r$ supplies all
corners and monochromatic edge positions.}
\label{tab:cuboid-demands}
\end{table}

A majority bucket with $4a$ sources supplies $Q$: use eight for its
corners by Lemma~\ref{lem:bucket}, and the remaining $4a-8$ for its
monochromatic edges.
For $H$, the corner multiset is
\[
 2uuv+2uuw+2uvv+uvw+uwv.
\]
It suffices to have $2a$ majority $T_u$ sources, $b$ majority $T_v$ sources,
a minority source supporting $\{u,v,w\}$, and one further majority source.
Four $T_u$ copies supply the $uuv,uuw$ corners and $2(a-2)$ fill the $uu$
edges; two $T_v$ copies supply $uvv$ and $b-2$ fill the $vv$ edges.
The minority and extra majority copies supply the two trichromatic
orientations. The $T_u$, $T_v$ and minority supplies are disjoint; the
extra majority copy is chosen outside the reservations.

For either row, the largest bichromatic demand plus $r+3$ is
$4a+2b+2c-5$. We have
\begin{equation}\label{eq:cuboid-budget}
 N\ge4a+2b+2c-5:
\end{equation}
if $b+c\ge6$ this follows from $N\ge L$, since the difference is
$2b+2c-11\ge1$; the remaining shapes have $a=b=2$, $c\le3$, and
$N\ge29$ suffices. Total demand is $L\le N$. Thus both constructions
extend to frames whenever the reserved sources are available.

Suppose no majority bucket has size $4a$. Put $h=4a-1$, and write
$t_u,p_u$ for the $T$ count and $P$ incidence in bucket $u$.
With $n_T,n_P$ the majority totals, double counting gives
\begin{equation}\label{eq:cuboid-buckets}
\begin{gathered}
 t_u+p_u\le h,\qquad n_T+2n_P\le4h,\qquad A=n_T+n_P\le4h,\\
 N\le16a+2,\qquad n_T\ge2A-4h\ge16a-14.
\end{gathered}
\end{equation}
There is a minority source, since otherwise $N=A\le4h$ contradicts
$N\ge16a-3=4h+1$. Every minority source supports a triple containing any
prescribed colour. An ordinary minority source supplies the opposite
orientation; a checkerboard supplies both.

\subsection{Completing the upper bound}
Since $L\le N\le16a+2$, integrality gives
\begin{equation}\label{eq:cuboid-width}
 b+c\le3a+4,\qquad 2b\le3a+4.
\end{equation}
Some $t_u\ge2a$, because $n_T\ge16a-14>4(2a-1)$ for every $a\ge2$.
If at most one colour has fewer than $b$ copies of its $T$ variety, a
minority triple containing $u$ has another colour $v$ with $t_v\ge b$.
These provide the $T_u$ and $T_v$ sources for $H$.

Otherwise two colours have fewer than $b$ copies, giving
\[
 n_T\le2(b-1)+2h=8a+2b-4\le11a.
\]
For $a\ge3$ this contradicts $n_T\ge16a-14$, so necessarily $a=2$.
Now $h=7$, $n_T\ge18$ and $L=4(b+c)-8\ge8b-8$. Hence
\[
 16b-56\le2L-40\le2N-40\le n_T\le2b+12,
\]
which forces $b\le4$. Thus $n_T\le20$, while
$n_T\ge2(N-M)-28\ge30-2M$ forces $M\ge5$.
Since every $t_i\le7$, at least two colours $u,v$ have $t_u,t_v\ge4$;
otherwise $n_T\le7+3\cdot3=16$. They satisfy $t_u\ge2a$ and $t_v\ge b$.
An ordinary minority source supports a triple containing both colours.
If all minority sources are checkerboards, at most three miss the pair
$uv$, by~\eqref{eq:cuboid-minority}. Since $M\ge5$, another checkerboard
supports a triple containing $u,v$: one missing $ij$ supports precisely
the triples omitting $i$ or $j$. Again the $H$ reservation is available.

In both cases, $A\ge16a-9>2a+b$, by~\eqref{eq:cuboid-width}, supplies the
extra majority copy. Equation~\eqref{eq:cuboid-budget} and the residual
Hall matching complete the frame. Together with
Corollary~\ref{cor:cuboid-D} and the three lower bounds, this proves the
frame formula in Theorem~\ref{thm:cuboids}.

Every position outside the frame exposes at most one face when all sides
are at least two. Every unused source contains every palette colour,
so it can fill any such position. Therefore $G=\max\{abc,F\}$.
Writing $a=x+2$, $b=y+2$, $c=z+2$ gives
\[
 abc-L=xyz+2(xy+xz+yz)\ge0,
\]
proving the complete-box formula and finishing Theorem~\ref{thm:cuboids}.
A side of length one would let a source expose opposite exterior faces
simultaneously and is outside this completion argument.

\subsection{Consequences}
The cube thresholds follow by setting $a=b=c=n$:
$F(n,n,n)=16n-3$ and $G(n,n,n)=\max\{n^3,16n-3\}$.
The cube stability theorem additionally classifies failures in its stated
range. The two thinnest families also follow immediately.

\begin{corollary}[Rods and strips]\label{thm:thin}
Put $K_\ell=\max\{29,\lfloor(24\ell-1)/5\rfloor\}$. Then
\[
 F(2,2,\ell)=G(2,2,\ell)=K_\ell \qquad(\ell\ge2),
\]
and
\[
 F(2,3,\ell)=K_\ell,\qquad G(2,3,\ell)=\max\{29,6\ell\}
 \qquad(\ell\ge3).
\]
\end{corollary}
Despite their four extra frame positions, strips have the same frame
thresholds as rods. Shape matters even at fixed cardinality:
$2\times2\times7$ and $2\times3\times6$ both have twenty-eight frame
positions, but their thresholds are thirty-three and twenty-nine.

\begin{corollary}[No spare sources]\label{cor:cuboid-volume}
A frame has threshold equal to its cardinality if and only if
\[
 b+c\ge3a+4\qquad\hbox{and}\qquad c\le5(a+b)-19.
\]
Every complete cuboid can be assembled from any collection of exactly
its volume, except for all rods $2\times2\times c$ and the four shapes
\[
 2\times3\times3,\quad2\times3\times4,\quad
 3\times3\times3,\quad3\times3\times4.
\]
The latter thresholds are $29,29,45,45$, respectively.
\end{corollary}
\begin{proof}
The frame conditions are precisely $L\ge16a-3$ and $L\ge K(c)$ after
integer rounding. For the second,
$24c-1<5(L+1)$ is equivalent to $c\le5(a+b)-19$.
Rods satisfy $K(c)>4c$. For every other ordered shape, $ab\ge6$, so
$abc\ge6c>K(c)$. The only remaining volume deficit is $abc<16a-3$.
For $a\ge4$, already $a^3\ge16a-3$. For $a=3$ it leaves
$(b,c)=(3,3),(3,4)$; for $a=2$, $b\ge3$, it leaves the same two pairs.
\end{proof}

For example, the next family after strips satisfies
\[
 F(2,4,c)=
 \begin{cases}
 29,&c=4,5,\\
 4c+8,&6\le c\le11,\\
 \lfloor(24c-1)/5\rfloor,&c\ge12.
 \end{cases}
\]
For $c\ge4$, any collection of $8c$ source cubes assembles a complete
$2\times4\times c$ box.
